\documentclass[10pt,twocolumn]{article}
\usepackage[scaled]{helvet}
\renewcommand\familydefault{\sfdefault}
\usepackage[T1]{fontenc}
\usepackage[margin=0.7in]{geometry}
\usepackage{titlesec}
\usepackage{enumitem}
\usepackage{parskip}
\setlength{\columnsep}{0.24in}
\setlist[itemize]{leftmargin=1.2em, topsep=0.2em, itemsep=0.18em}
\titleformat{\section}{\large\bfseries}{\thesection}{1em}{}

\begin{document}
\title{\vspace{-1.6cm}MLflow on AWS}
\author{AWS ML Study Notes}
\date{}
\maketitle
\small

\section*{Overview}
MLflow provides experiment tracking, model registry style capabilities, and artifact management for machine learning work. Running it on AWS makes it accessible to a team while using AWS storage, databases, and security controls.

Its real value is organizational memory. It records what ran, with which parameters, on which data, and with what results so model development becomes auditable and comparable.

\section*{Core Concepts}
\begin{itemize}
\item A typical MLflow setup includes a tracking server, a backend store for metadata, and an artifact store such as S3 for model files and plots.
\item Runs capture parameters, metrics, tags, and artifacts. Registered models organize versioned outputs that teams may later evaluate or deploy.
\item MLflow does not replace the entire production platform. It complements it by handling experiment lineage and model packaging at the experimentation layer.
\end{itemize}

\section*{How It Works}
\begin{itemize}
\item A notebook, script, or training job logs metrics and artifacts to the tracking server while the run is executing.
\item Team members compare runs, inspect outputs, and decide which model versions deserve promotion, registration, or deeper evaluation.
\item When integrated with an MLOps pipeline, an approved run can trigger packaging, validation, or deployment workflows in AWS.
\end{itemize}

\section*{AI and ML Relevance}
\begin{itemize}
\item MLflow is useful when multiple experiments, datasets, and parameter combinations are being tested and human memory is no longer reliable enough to track what worked.
\item It also provides a bridge between ad hoc experimentation and operational model release by giving each promising artifact a stable identity and metadata trail.
\end{itemize}

\section*{Operational Notes}
\begin{itemize}
\item Protect the tracking server with authentication, network isolation, and backups for the metadata store. Experiment history is operationally important.
\item Separate artifact storage from code repositories. Model outputs can be large and need lifecycle and retention management.
\item Decide early which metadata fields are mandatory, or the run history will become inconsistent and hard to search.
\end{itemize}

\section*{What To Remember}
\begin{itemize}
\item MLflow answers what happened during experimentation; deployment systems answer what is live in production.
\item Without disciplined logging, MLflow becomes a dump of half labeled runs. Governance needs conventions.
\item For teams iterating quickly, experiment tracking is not optional once more than one person or one model is involved.
\end{itemize}

\end{document}
